\chapter{Introduction}
\begin{quote}
  \itshape A \emph{Sefer Torah} (in Hebrew \texthebrew{ספר תורה},
  literally ``Book of Torah'') is a handwritten copy of the Pentateuch
  in the form of a scroll, used in ritual Torah reading during Jewish
  prayers. The rules concerning the composition of Torah scrolls are
  strict: a deficiency in parchment processing, ink properties or a
  transcription error render the scroll invalid for ritual use.
\end{quote}

\begin{multicols*}{2}
  Sefer is a purely functional data-processing language. It provides
  static, polymorphic typing, algebraic datatypes, typeclasses,
  pattern-matching and a concise yet powerful syntax adapted for both
  one-liners and more complex scripts. At the core of every Sefer
  program lies the \verb|Iterator| typeclass and the many iterator
  adapters defined in its standard library.

  Sefer being a purely functional language, no input/output (I/O)
  access is provided to the user. The only way for a Sefer program to
  interact with the outside world is to read input from an iterator
  over the characters (or bytes when binary mode is used) written on
  the standard input and to return an expression, which will
  afterwards be displayed on the standard output by the virtual
  machine when it terminates.

  \section{Rationale}

  Sefer's main reason for existence is to replace AWK as the go-to
  utlity for data processing on UNIX-like systems. There are several
  motivations for this, which are all linked to each other. AWK was
  designed in the 1970s and its most widely used form (GNU AWK) was
  published in the late 1980s; AWK is tragically a product of its
  time, in both its internal and external design. At that time, Pascal
  and Fortran were mainstream languages and the ANSI C standard was
  just being finalised. Actually, when AWK was created in 1977, Brian
  Kernighan (who is also the `K' in ``AWK'') and Denis Ritchie's
  seminal book, \emph{The C Programming Language} (1978), had not even
  been published yet.

  An AWK program is a series of pattern/action pairs, written as
  \verb|pattern { action }|. The input is split into \emph{records}
  (by default lines) and the \verb|action| is executed on every record
  that matches the \verb|pattern|. Two special patterns, namely
  \verb|BEGIN| and \verb|END| specify actions to be executed before
  and after the program, respectively. Inside actions, every record
  can be split into fields (on a given separator) which are accessible
  through the \verb|$n| variables, where \verb|n| is a number. This
  means that the entiry of AWK is as powerful as the Sefer combinator
  \verb|map_cond|, which maps iterator elements using a different
  function depending on the value of a predicate (see the standard
  library specification for more information on this adapter). Such
  lack of expressivity is not too much of a problem on simple
  one-liners, but quickly grows to become one as task complexity
  increases.

  Even on one liners, AWK suffers from a lack of clarity due to its
  imperative nature; indeed, declarative programming is especially
  well-suited for data processing, as the user extracting certain
  information from the data is more interested in the \emph{result} of
  the extraction than in its \emph{process}. See for example the
  following AWK program, which displays the number of words in the
  input (note that \verb|NF| holds the number of fields in the record):
  \begin{lstlisting}[language=Awk,caption={\texttt{wc -w} as an AWK
        program}]
      {
        words += NF
      }
      END { print words }
    \end{lstlisting}

  In Sefer, a similar program could be written as follows:
  \begin{lstlisting}[caption={\texttt{wc -w} as a
        Sefer program}]
    ||- (words | count)
     | sum
  \end{lstlisting}

  \verb|words| simply constructs an iterator from a string (which is
  actually but an iterator on characters), splitting it on every
  whitespace, \verb|count| consumes an iterator and returns the number
  of elements it yielded and \verb|sum| consumes an iterator,
  returning the sum of its elements. Another advantage of the
  declarative paradigm as imposed by Sefer is that the language itself
  can optimise computation in the backend, while guaranteeing a
  clear and concise frontend.

  AWK is not a bad tool in and of itself, rather a tool that was made
  for low-resource computers in an era in which the imperative
  paradigm was quasi-ubiquitous outside of theoretical
  experiments. AWK is an out-of-date tool that needs a modern
  replacement and that is what Sefer aims to be.

\end{multicols*}
